\section{FISH}
\label{sec:method}
FISH is a toolset with four parts: \emph{instrumentation}, tracepoints added to the ROS~2 client libraries plus the CUDA profiler; \emph{acquisition}, an entry point and a daemon that run an unmodified application under that instrumentation; \emph{extraction}, which turns the trace into the model of Fig.~\ref{fig:model}; and \emph{FISH-eye}, a viewer that renders the model at any granularity the expansion operator allows (Fig.~\ref{fig:screens}). FISH currently targets ROS~2 Humble.

\noindent\textbf{Application structure.} FISH builds on ros2\_tracing, whose initialisation events record each object as it is created and whose execution events record what then happens to it, linked by the object's handle. To its 14 events FISH adds 28, in five categories that close the gaps of Section~\ref{sec:goal}: finalize events, which end the lifetime the stock events begin, so a reused handle is not mistaken for its predecessor; introspection events for each executor's type and thread count and each callback's group; start and end events for clients and waitables, the intra-process subscription waitable among them, which complete the five kinds the executor dispatches, each recorded with the thread it ran on; link events at the publish, receive and request sites, which tie a message to the callback that issued it; and registration events that give rclpy nodes the same chain as rclcpp. Extraction lays down the objects and the declared edges from the initialisation events (Fig.~\ref{fig:model}a), decides the remaining edges from the trace patterns of Fig.~\ref{fig:model}c, and lifts every edge along the containment (Fig.~\ref{fig:model}b). The result is checked against what the traced binaries create: a node constructor must yield one node, seven entities and seven callbacks, and the model is complete only if every such expectation appears in it.

\noindent\textbf{CPU--GPU interplay.} GPU work is recorded by CUPTI through \texttt{nsys}, aligned to the LTTng clock at ingest, bound through correlation ids to the API call and through the calling thread's active callback to the callback that issued it. The scheduler's context switches split each callback execution into on-core, preempted and blocked time, and the GPU timeline shows which of it was spent waiting for submitted work. CUPTI sees only processes started under the profiler, and a container cannot be restarted without losing the components loaded into it, so profiling has to be attached at launch: the FISH entry point wraps \texttt{ros2 launch}, reads the launch description and spawns the containers and nodes that will use the GPU under \texttt{nsys} from the outset, while the daemon polls \texttt{/proc} and restarts under the profiler any process outside the launch tree that holds a CUDA context. Per invocation, the recorded kernels, copies and synchronisations become the typed DAG of Fig.~\ref{fig:model}d; invocations of one signature form a shape, and the steady-state shapes fold into one conditional graph.

\noindent\textbf{Task chains.} The daemon gates the run with signal files: it polls the component list and the number of started processes, and only once both have been stable over several polls does it mark the system stable, complete the GPU hand-over and start the input, a rosbag replay or a live source. Phase boundaries then come from the trace itself: initialisation ends when every periodic timer has fired once, warm-up when every recurring callback has completed its first execution, and a callback that never ran is kept as such. A chain is read off the model as a path at the callback layer, its GPU segments one expansion below it and its nodes, executors and processes any number of lifts above.
